第二章已经表明，对于由物理本构与显式几何约束主导的病态系统，可以通过重构顺应内在力学规律的广义坐标来隔离高频刚度。然而，在更一般的三维有限元连续体分析场景中，病态根源往往来自空间离散本身，而不再表现为显式的宏观物理约束。特别是当网格包含高度畸变的单元结构时，局部基函数的数值性质会急剧恶化，并通过高频人造刚度干扰全局低频响应。本章聚焦这一由离散化诱发的软硬耦合问题。

\section{研究背景}

在基于有限元的底层弹性物理模拟中，系统偏微分方程数值求解的稳定性，在很大程度上受制于最底层支撑网格的几何质量。与细杆问题中源自物理本构的极短距离约束不同，本章所处理的高频刚度成分，隐藏在局部拉格朗日基函数内部的梯度变化之中。当网格区域产生诸如严重扁平化、扭曲变斜或者具有一对非常靠近的强连接顶点（如典型的“薄片”单元，Slivers）时，由参考规整单元向物理真实单元映射的雅可比矩阵将趋于奇异；这会导致形函数空间梯度显著增大，进而在局部刚度矩阵中引入较高的高频数值刚度。这些局部高频响应可能推高全局刚度矩阵的条件数上限并引发“锁定”现象（\Cref{fig:locking_issue}），导致克雷洛夫等迭代求解器的收敛速度变慢。

\subsection{相关工作与传统方法的局限}

为应对此类畸变诱发的数值问题，现有方法大致可分为三类。第一类从几何层面改善网格质量，通过重网格化、拓扑修复或节点平滑等手段尽量消除劣质单元。已有研究表明，单元形状质量与离散误差、刚度矩阵条件数之间存在紧密联系~\cite{Shewchuk2002whatIsAGoodLinearFiniteElement, Shewchuk2007Anisotropic, Tournois2009PerturbingSlivers, Cheng1999SliverExudation, RH2005EliminateSlivers}。但在复杂边界、频繁切割与拓扑变化的动态场景中，这类方法往往成本较高，也难以严格避免畸变单元的持续出现。

第二类方法从代数层面缓解病态性，例如多尺度不完全 Cholesky 分解、代数多重网格（AMG）预处理以及多层 Schwarz 算子等~\cite{chen_multiscale_2021, Griebel2003amgElas, Wu2022schwarz}。它们能够改善迭代收敛，但其粗层空间仍建立在底层畸变基函数之上，因此对由局部梯度奇异所诱发的高频人造刚度缺乏直接穿透能力。

第三类工作则将关注点回到离散基函数本身，包括数值粗化与同质化、算子自适应基函数、多面体有限元与虚单元法等~\cite{Geers2010, Lukas2020, Bacigalupo2014, OstojaStarzewski2006, ChenDdFEM2015, li2023NeuralMetamaterialNetworks, chen2018numerical, heFast2023, chen2019material, Li2022CBN, ni_numerical_2023, SukumarMalsch2006, Jabbari2022AdaptiveRefinement, Dasgupta2003, Chen2020PrismsCones, BeiraoDaVeiga2013a_VEMBasic, BeiraoDaVeiga2013b_VEMHitchhiker, GainesJenkins2015}。这类研究表明，通过重构局部解空间而非单纯修补几何或增强代数预处理，有可能更直接地削弱病态单元对全局求解的干扰。本文正是在这一思路下，面向畸变网格局部区域构造算子感知的基空间重构与隔离机制。

\subsection{本章方法与贡献}

针对上述问题，本章提出了一种基于算子感知的局部基空间重构与隔离方法，旨在缓解由离散化导致的病态性，并降低局部高频人造刚度对全局低频响应的数值干扰。我们不再采用全面重构三维体全局自由度的做法，而是采用局部策略：首先利用基于残差的自适应细化（h-refinement）技术定位隐藏的低质量强连接与退化区间；继而将这些区域局部聚合为多面体，并在聚合后的微区构造算子自适应驱动的矩阵值基函数。由此形成局部层次结构，并将原本易激发高频刚度的离散表征重构为更稳定的局部基空间。

对于包含强连接构型的单元（由两个间距极短节点造成的近乎奇异性约束），本章进一步引入了 Dirac-wavelet 技术。该技术将强耦合的数值方向从主求解空间中显式隔离，削弱了不良节点对全局系统特征谱最大特征值的影响。整个重构与隔离流程在降低系统矩阵条件数的同时，确保了新基函数在计算域中维持紧凑支撑形态，可与常规求解器配合使用。

本文的主要贡献总结如下：
\begin{itemize}
    \item 提出了一套面向畸变网格局部问题区域的算子感知基空间重构框架。该框架通过对坏形状单元形成的局部 patch 进行聚合，并在聚合后的区域内构造局域化基函数，在基本不改变原始自由度规模的前提下，缓解由离散退化诱发的人工高频刚度；同时配套给出了一种排序与判定策略，用于识别并选择需要聚合处理的单元区域。
    \item 提出了一种针对强连接刚性模态的显式隔离机制，并通过数值实验测试了其适用性。对于由强连接或近奇异局部结构引起的特殊高频硬模态，本文借助 Dirac-wavelet 构造将其从主求解空间中识别与处理；结合基准测试与应用场景分析，评估了该局部重构与隔离策略在改善谱性质和求解效率方面的效果。
\end{itemize}
